\section{Open-Source Software}

\begin{itemize}[]
    \item \href{https://github.com/salehisaeed/TensorforceFoam}{\texttt{TensorforceFoam}}\swref{Salehi2024TensorforceFoam}: A TensorFlow-based intrusive Deep Reinforcement Learning (DRL) framework for OpenFOAM. The Tensorforce package is utilized for the DRL computations. The term intrusive refers to the direct integration of the DRL agent within the CFD environment. Such integration eliminates the need for any external information exchange during DRL episodes. The framework is parallelized using the message passing interface (MPI) for Python (mpi4py) to manage parallel environments for computationally intensive CFD cases through distributed computing.

    \item \href{https://github.com/salehisaeed/agentJet}{\texttt{agentJet}}\swref{Salehi2024agentJet}: An OpenFOAM boundary condition that deploys a trained Deep Reinforcement Learning (DRL) policy as a jet actuator directly inside the CFD solver. The saved TensorFlow policy is loaded and evaluated through the TensorFlow C++ API (cppflow), so no Python coupling or external communication is needed at run time. At each control period, the flow state is sampled at user-defined probe locations and normalized, the policy (deterministic or stochastic) returns a jet velocity, which is ramped up smoothly from the previous action, and the state and action trajectory is written to file for analysis.

    \item \href{https://github.com/salehisaeed/semiImplicitSlip}{\texttt{semiImplicitSlip}}\swref{Salehi2022semiImplicitSlip}: A mesh motion library in OpenFOAM that implements a robust semi-implicit algorithm for slipping the mesh points on curved surfaces. The algorithm includes two steps, namely, an explicit step based on the general slip condition and an implicit step based on the Dirichlet condition. It employs the Laplacian smoothing equations to spread the mesh deformation into the domain. Additionally, a solid-body rotation may be added on top of the deformed mesh, which could be useful for modeling the runner region in transient operation of Kaplan turbines which contains simultaneous mesh deformation and solid-body rotation of the mesh.

    \item \href{https://github.com/salehisaeed/OpenFOAM-Swirl-Number}{\texttt{OpenFOAM-Swirl-Number}}\swref{Salehi2025SwirlNumber}: A complete pipeline for computing the swirl number, the ratio of the axial flux of angular momentum to the axial flux of axial momentum, in OpenFOAM. A set of modular function objects extracts the velocity components, transforms them to cylindrical coordinates, and integrates the tangential and axial momentum fluxes over an arbitrary sampled cross-section, and a Python script computes and plots the swirl number over time. It is demonstrated on the \texttt{axialTurbine\_rotating\_oneBlade} tutorial case.

    \item \href{https://dataverse.harvard.edu/dataset.xhtml?persistentId=doi:10.7910/DVN/31JGOM}{\texttt{Francis-99 transients}}\swref{Salehi2021Francis99Transients}: Open-source case and code for CFD simulation of the transient operation of a Francis turbine in OpenFOAM.

    \item \href{https://github.com/salehisaeed/OSCFD}{\texttt{OSCFD}}\swref{Salehi2024OSCFD}: Test codes and examples for the PhD course CFD with Open-Source Software at Chalmers University of Technology, complementing the lectures on memory management in C++ and OpenFOAM. Standalone C++ examples cover dynamic allocation, destructors, shallow and deep copies, copy constructors, assignment overloading, and smart pointers, while dedicated OpenFOAM applications examine its own mechanisms, \texttt{autoPtr} and \texttt{tmp}, for fields and matrices, including a memory-test variant of \texttt{simpleFoam}.
\end{itemize}
